\chapter{Stress}

\section{Introduction}
{\bf Stress} can generally be thought of as a ratio
of force to area.  The greater the force for a given area, the higher
the stress.  The greater the area for a given force, the lower the
stress.  

Forces and areas are vector quantifies, 
with magnitude and direction.  

The force magnitude, a scalar value, must be zero or positive.  
Negative force values are disallowed because the force magnitude
is the length of the force vector, 
which is independent of force direction.

The area magnitude, a scalar value, must be positive.
A zero area value is disallowed because division by zero is
undefined.  Negative area values are disallowed because the 
area magnitude is the length of the area vector, 
which is independent of area direction.  

The force direction is described by a unit vector in a particular
reference frame.  The unit vector endows 
the force vector signs, negative and positive, to indicate
direction with the reference frame.  

Similarly, the area direction is described by a unit vector
in a particular reference frame.  This reference frame is typically
the same reference frame used to describe the force.  However, if
the two references are not the same, then a basis transformation may
be used to put both vectors into the same reference frame. 

\section{Stress Measures}
\begin{table}[htb]
\caption{Seth-Hill family of strains 
and their respective stress conjugates.}
\label{tab:table_of_conjugate_stress}
\begin{center}
\begin{tabular}{|c|c|c|c|} \hline
     & {\bf Name} & {\bf Name} \\ 
 $m$ & and & and \\ 
     & {\bf Strain Tensor} & {\bf Stress Tensor} \\ \hline \hline
 & Green-Lagrange & Piola-Kirchhoff Two \\
 2  
   & $\tEbold_{(2)} = \dst\frac{1}{2}\left(\tUbold^2 - \vid\right)$ 
   & $\tSbold = J \;\tFbold^{-1} \ve{\sigma} \tFbold\;\minustranspose$
   \\
   & & \\ \hline
 & engineering (Biot, nominal) & Biot (Jaumann) \\
 1  
   & $\tEbold_{(1)} = \tUbold - \vid$
   & $\tBbold = 
   \half \left(\tSbold \tUbold + \tUbold \tSbold \right)$
  \\
   & & \\ \hline
 & log (Hencky, natural) & \\
 0  
   & $\tEbold_{(0)} = \ln\tUbold$
   & to come
   \\
   & & \\ \hline
 & true & Swainger \\
 -1  
   & $\tEbold_{(-1)} = \vid - \tUbold^{-1}$
   & {\sc SW}$= \half J \left(\tUbold \hat{\ve{\sigma}} + 
           \hat{\ve{\sigma}} \tUbold \right)$
   \\
   & & \\ \hline
 & Euler-Almansi & Almansi\\
 -2  
   & $\tEbold_{(-2)} = \dst\frac{1}{2}\left(\vid - \tUbold^{-2}\right)$
   & $\tAbold = J \tFbold\transpose \ve{\sigma} \tFbold$
   \\
   & & \\ \hline
\end{tabular}
\end{center}
\end{table}

The {\bf Kirchhoff stress} $\ve{\tau}$ (also known as the 
weighted Cauchy stress) is defined as
\be
 \ve{\tau} \defe J \; \ve{\sigma}. 
\ee
The Kirchhoff stress is often used in metal plasticity formulations.

The {\bf Co-rotated Cauchy stress} (also known as the back-rotated
Cauchy stress) is defined as
\be
 \hat{\ve{\sigma}} \defe \tRbold\transpose \ve{\sigma} \; \tRbold
\ee

The {\bf Piola-Kirchhoff One stress $\tPbold$} is defined as
\be
 \tPbold \defe J \; \ve{\sigma}\; \tFbold\minustranspose 
 = \ve{\tau} \;\tFbold\minustranspose
\ee

The {\bf Piola-Kirchhoff Two stress $\tSbold$} is defined as
\be
 \tSbold
 \defe J \; \tFbold\inverse \;\ve{\sigma} \; \tFbold\minustranspose 
 = \tFbold\inverse \; \ve{\tau}\; \tFbold\minustranspose.
\ee
Thus
\be
 \tPbold = \tFbold \tSbold.
\ee


\section{Stress Rates}
Stress rates arise when the stress depends on variables that have
time dependency.  Viscoelastic material behavior, for example, will
exhibit strains, and thus stresses, with time dependence.  

